Two metal--ferroelectric--metal capacitors were compared. Stack~A was TaN/\hzo{}(10~nm)/TiN on Si. Stack~B was TiN/PE-TiN/\hzo{}/WO$_x$/W on Si, where the deliberate WO$_x$ layer modifies the bottom interface and provides a second stack with a different interfacial impedance. Grazing-incidence XRD was used to confirm crystallization of the ferroelectric films, while the electrical study focused on the relation between the two remanent branches of the small-signal dielectric response.

Frequency-resolved $C$--$V$ measurements were performed in parallel $C_p$--$R_p$ mode with a 50~mV rms signal at 61 logarithmically spaced frequencies from 1~kHz to 4~MHz. The dc bias was swept $-2 \rightarrow +2 \rightarrow -2$~V so that the two polarization branches were collected under identical conditions. For each frequency, the zero-bias capacitive memory window was obtained from the branch separation,
\begin{equation}
\CMW_\xi(f)=\xi_\downarrow(0,f)-\xi_\uparrow(0,f),
\label{eq:cmw_short}
\end{equation}
where $\xi_\uparrow$ and $\xi_\downarrow$ denote the branches reached after opposite preset states. The crossing field $\Ecross$ was taken from the bias at which the two branches became capacitively indistinguishable.

To test whether changes in the window simply mirrored changes in polarization, the capacitive data were paired with PUND and switching-current measurements. The same devices were measured at 300~K and 10~K so that the zero-bias readout and the switching response could be compared on the same temperature axis. High-frequency interpretation was limited to the range below the series-resistance-dominated region identified from the measured roll-off and standard dielectric-analysis practice~\cite{nicollian1982,schroder2006}.
